\ifdefined\CRevisionRound\else\def\CRevisionRound{2}\fi
\documentclass[10pt]{article}
\pdfvariable suppressoptionalinfo 767
\usepackage[a4paper,margin=19mm]{geometry}
\usepackage{amsmath,amssymb,amsthm,mathtools,booktabs,array}
\usepackage{fontspec}
\setmainfont{TeX Gyre Pagella}
\setsansfont{TeX Gyre Heros}
\setmonofont{Latin Modern Mono}
\newfontfamily\cjkfont{Droid Sans Fallback}
\usepackage[hidelinks]{hyperref}
\newcommand{\RR}{\mathbb R}
\newcommand{\ZZ}{\mathbb Z}
\newcommand{\TT}{\mathbb T}
\newcommand{\e}{\mathrm e}
\newcommand{\ii}{\mathrm i}
\newcommand{\dd}{\,\mathrm d}
\newcommand{\Tr}{\operatorname{Tr}}
\newcommand{\Fix}{\operatorname{Fix}}
\newcommand{\sgn}{\operatorname{sgn}}
\newtheorem{theorem}{Theorem}
\newtheorem{lemma}{Lemma}
\ifcase\CRevisionRound
\title{A Charged Particle on a Flat Magnetic Torus:
Clean Cyclotron Return and Sharp Flux Quantization}
\or
\title{A Charged Particle on a Flat Magnetic Torus:
Clean Return, Flux Quantization, and Signed Landau Multiplets}
\else
\title{A Charged Particle on a Flat Magnetic Torus:
Clean Cyclotron Return, Landau Multiplets, and Two Revival Times}
\fi
\author{HCS-C376 theorem package}
\date{4 September 2026}

\begin{document}
\emergencystretch=3em
\maketitle

\begin{abstract}
\ifcase\CRevisionRound
We close, in one normalization, the classical and quantum dynamics of a
charged unit-mass particle on a rectangular flat two-torus with nonzero
constant magnetic field.  Every positive-energy classical trajectory has
least period $2\pi/|B|$; each return is maximally clean, whereas no other
time fixes a point of an energy shell.  The cyclotron center identifies the
orbit quotient as a two-torus.  A quantum line bundle with curvature
$-\ii B\dd x\wedge\dd y$ exists precisely at integral flux
$N=BL_xL_y/(2\pi)$.  Necessity follows from Chern--Weil theory and
sufficiency from an explicit rectangle cocycle.  This first round closes
the global classical return and topological existence questions only.
\or
We fix one normalization for a charged unit-mass particle on a rectangular
flat two-torus.  At nonzero field, every positive-energy classical orbit has
least period $2\pi/|B|$, and every return is maximally clean.  Curvature
$-\ii B\dd x\wedge\dd y$ admits a Hermitian line bundle exactly when
$N=BL_xL_y/(2\pi)$ is integral.  On that bundle the physical Bochner
Hamiltonian has levels $|B|(n+1/2)$, each with multiplicity $|N|$ for every
flat holonomy.  For the fixed positively oriented division vectors, the
magnetic translations obey $UV=\zeta_{|N|}^{\sgn N}VU$; hence negative flux
carries the conjugate clock--shift representation rather than an erased
sign.  This round proves the clean return, global quantization, complete
Landau ladder, and signed finite-Heisenberg action.
\else
We close, in one normalization, the classical and quantum dynamics of a
charged unit-mass particle on a rectangular flat two-torus with nonzero
constant magnetic field.  Every positive-energy classical trajectory has
least period $2\pi/|B|$, and each return is maximally clean.  Quantum
curvature $-\ii B\dd x\wedge\dd y$ exists globally precisely at integral
flux $N=BL_xL_y/(2\pi)$.  The physical Bochner Hamiltonian has levels
$|B|(n+1/2)$ of multiplicity $|N|$; fixed positive division vectors retain
the field orientation through $UV=\zeta_{|N|}^{\sgn N}VU$.  We derive the
heat trace, Hurwitz spectral zeta, scale-independent determinant
$\det_\zeta H_B=2^{|N|/2}$, and two distinct least times: one classical
period gives the scalar $-I$, whereas two give the identity.  The boundary
atlas includes the exact lattice-normalized zero-field closure criterion.
Topological flux integrality is not arithmetic provenance, so the strict
Route-A tuple begins with A0 fail and the candidate is rejected.
\fi
\end{abstract}

\begin{center}{\cjkfont 中文摘要}\end{center}
{\cjkfont\small\raggedright
\ifcase\CRevisionRound
本文在同一套符号和物理单位中研究矩形平坦二环面上的恒定磁场运动。
经典部分不从局部轨道猜测周期，而是在普适覆盖上直接积分方程：
动量作匀速旋转，回旋中心保持不变，位置由同一旋转精确重建。
由此证明任意正能量轨道的最小相空间周期均为二派除以磁场绝对值；
非返回时刻没有不动点，返回时刻整个能量层连同其切空间同时固定，
因此所得周期族是最大干净的连续族，其轨道商仍为二环面。
量子存在性则由全局拓扑决定。
对曲率负虚数乘以磁场面积形式的厄米线丛，第一陈数恰为磁通除以二派；
只有该数为整数时对象存在。必要性来自陈类积分，
充分性由基本矩形上的显式过渡函数给出。
本轮只封闭经典回归与磁通量子化，不预告后续谱结论。
\or
本文统一处理平坦磁二环面的经典回归、全局量子化和朗道能级。
经典方程在普适覆盖上可完全积分，守恒回旋中心把所有正能量轨道锁定到共同的最小周期，
并说明返回映射在整个能量层上都是恒等，从而得到最大干净而非孤立的周期族。
曲率积分给出有符号磁通整数，它既是线丛存在的充要条件，也是线丛次数；
任意平坦荷洛诺米只改变连接而不改变能谱。
采用带零点半位移的正博赫纳哈密顿量后，
升降算符与亏格一的黎曼罗赫定理共同给出完整等距能级及每层磁通绝对值重简并。
为避免抹去负磁通，本文固定正向的两条细分格矢，
证明磁平移交换相位为单位根的磁通符号次幂，
并在每个能级上构造相应的共轭钟移表示。
本轮止于完整朗道阶梯和有符号有限海森堡作用。
\else
本文在固定物理归一化下封闭恒定磁场中矩形平坦二环面的经典与量子动力学。
经典解由刚性动量旋转和守恒回旋中心组成，
因此每条正能量轨道都有共同最小周期，且返回映射在整个能量层上为恒等；
这是一族最大干净轨道而非孤立周期点。
量子对象存在当且仅当总磁通给出整数第一陈数。
正博赫纳哈密顿量的升降代数和环面黎曼罗赫定理给出带零点半位移的全部朗道能级及磁通绝对值重简并。
固定正向细分格矢后，磁平移交换相位保留磁通符号，负磁场对应共轭钟移表示。
由完整谱进一步求得热迹、赫维茨谱泽塔、与磁场尺度无关的泽塔行列式，
并严格区分一个经典周期产生的负恒等标量复现和两个周期产生的真正恒等复现。
边界分析还给出零场直线流按环面格矢归一化的闭轨条件，
分别处理轴向、静止与稠密情形。
最后指出陈数整数性属于拓扑而非素数算术，故严格路线甲判据在首道算术门即失败。
\fi
\par}

\medskip
\ifcase\CRevisionRound
\noindent\textbf{Keywords:} magnetic torus; cyclotron flow; clean return; flux quantization; line bundle

{\cjkfont\noindent 关键词：磁环面；回旋流；干净回归；磁通量子化；线丛}
\or
\noindent\textbf{Keywords:} magnetic torus; Bochner Laplacian; Landau levels; flux quantization; magnetic translations; finite Heisenberg group

{\cjkfont\noindent 关键词：磁环面；博赫纳拉普拉斯；朗道能级；磁通量子化；磁平移；有限海森堡群}
\else
\noindent\textbf{Keywords:} magnetic torus; Landau levels; magnetic translations; spectral zeta; zeta determinant; quantum revival; Route A

{\cjkfont\noindent 关键词：磁环面；朗道能级；磁平移；谱泽塔；泽塔行列式；量子复现；路线甲}
\fi

\section{Classical magnetic flow}

Let
\[
 \TT^2=\RR^2/(L_x\ZZ\oplus L_y\ZZ),\qquad A=L_xL_y,
\]
where $L_x,L_y>0$.  On $T^*\TT^2$ use lifted coordinates
$(q,p)=(x,y,p_x,p_y)$ and the twisted form and Hamiltonian
\begin{equation}
 \omega_B=\dd x\wedge\dd p_x+\dd y\wedge\dd p_y
       +B\dd x\wedge\dd y,
 \qquad H_{\rm cl}=\tfrac12|p|^2.                 \label{eq:classical}
\end{equation}
Fix $B\ne0$ and write $J(u,v)=(-v,u)$.  Our Hamiltonian convention gives
\begin{equation}
 \dot q=p,\qquad \dot p=BJp.                     \label{eq:lorentz}
\end{equation}

\begin{theorem}[Complete return and clean fixed sets]\label{thm:classical}
On every energy shell $\Sigma_E=\{H_{\rm cl}=E\}$ with $E>0$, every orbit
has the same least period
\[
 T_{\rm cl}=\frac{2\pi}{|B|}.
\]
For $t\notin T_{\rm cl}\ZZ$, $\Fix(\Phi_t|_{\Sigma_E})$ is empty.  For
$t\in T_{\rm cl}\ZZ$, the return is the identity and
\[
 \Fix(\Phi_t|_{\Sigma_E})=\Sigma_E,\qquad
 \ker(D\Phi_t-I)=T\Sigma_E.
\]
Thus every nonzero return is maximally clean.  The orbit space of the
circle action on $\Sigma_E\cong\TT^2\times S^1$ is a two-torus.
\end{theorem}

\begin{proof}
Let $R_\theta=\exp(\theta J)$.  Direct integration of
\eqref{eq:lorentz} gives
\begin{equation}
 p(t)=R_{Bt}p_0,\qquad
 c=q+\frac{Jp}{B}=q_0+\frac{Jp_0}{B},\qquad
 q(t)=c-\frac{J R_{Bt}p_0}{B}.                   \label{eq:solution}
\end{equation}
This formula is well defined after reducing $q$ modulo the lattice and
proves completeness.  Since $E>0$, $p_0\ne0$.  A phase-space return forces
$R_{Bt}p_0=p_0$, hence $Bt\in2\pi\ZZ$; the position then returns as well.
This proves both the period and its minimality.  At an integral multiple of
$T_{\rm cl}$, formula \eqref{eq:solution} is the identity for all initial
data, so its derivative is the identity.  At every other time its rotation
has no nonzero fixed momentum.  These two observations give the stated
fixed sets and clean tangent equality.  Finally $c$ modulo the lattice
labels an orbit, proving the quotient assertion.
\end{proof}

\paragraph{Geometric meaning.}
The continuous center $c$ is the classical accidental symmetry.  Compact
configuration does not isolate trajectories: at a return time the entire
three-dimensional energy shell is fixed.  This fact is decisive for the
Route-A verdict below.

\section{Flux integrality and the quantum object}

We use units $\hbar=q_{\rm charge}=m=1$.  A quantum magnetic field means a
Hermitian line bundle $\mathcal L\to\TT^2$ with unitary connection $\nabla$
whose curvature is
\begin{equation}
 F_\nabla=-\ii B\dd x\wedge\dd y.                \label{eq:curvature}
\end{equation}

\begin{theorem}[Sharp flux condition]\label{thm:flux}
A pair $(\mathcal L,\nabla)$ satisfying \eqref{eq:curvature} exists if and
only if
\begin{equation}
 N=\frac{BA}{2\pi}\in\ZZ.                        \label{eq:flux}
\end{equation}
Its degree is $N$.  For a fixed $N$, tensoring by a flat line bundle gives
the two holonomy parameters and does not change the curvature.
\end{theorem}

\begin{proof}
Chern--Weil theory gives
\[
 c_1(\mathcal L)[\TT^2]=\frac{\ii}{2\pi}
 \int_{\TT^2}F_\nabla=\frac{BA}{2\pi},
\]
which proves necessity.  Conversely, for $N\in\ZZ$ take on the fundamental
rectangle the gauge $A_y=Bx$, $A_x=0$.  Identify the vertical sides
trivially and identify $x=L_x$ with $x=0$ by the unitary multiplier
$\exp(\ii BL_xy)$.  Its change around the horizontal cycle is
$\exp(\ii BA)=\exp(2\pi\ii N)=1$, so the corner cocycle closes and defines
the required bundle and connection.  Multiplying the two gluing maps by
constant phases produces all flat holonomy twists.
\end{proof}

\noindent\textit{Round-zero certificate: round zero proves complete clean
classical return and the sharp global flux condition.}

\ifnum\CRevisionRound>0
\section{Bochner Landau ladder}

Put $\Pi_j=-\ii\nabla_j$.  Equations \eqref{eq:curvature} and
\eqref{eq:flux} give
\begin{equation}
 [\Pi_x,\Pi_y]=\ii B,
 \qquad H_B=\tfrac12(\Pi_x^2+\Pi_y^2)=\tfrac12\nabla^*\nabla. \label{eq:HB}
\end{equation}
This physical Bochner normalization is important: it includes the
zero-point half-shift and is not the unshifted Dolbeault number operator.

\begin{theorem}[Spectrum, multiplicity, and holonomy rigidity]
Assume $N\ne0$.  The self-adjoint elliptic operator \eqref{eq:HB} has
\begin{equation}
 \operatorname{Spec}(H_B)=
 \{E_n=|B|(n+\tfrac12):n\in\ZZ_{\ge0}\}.         \label{eq:levels}
\end{equation}
Every $E_n$ has multiplicity $|N|$.  The level locations and
multiplicities are independent of both flat holonomies.
\end{theorem}

\begin{proof}
It is enough to take $B>0$; negative field is obtained by complex
conjugation and interchange of chirality.  Define
\[
 a=\frac{\Pi_x+\ii\Pi_y}{\sqrt{2B}},\qquad
 a^*=\frac{\Pi_x-\ii\Pi_y}{\sqrt{2B}}.
\]
Then $[a,a^*]=1$ and
\begin{equation}
 H_B=B(a^*a+\tfrac12),\qquad [H_B,a]=-Ba.         \label{eq:ladder}
\end{equation}
In the complex structure of the rectangular torus, $\ker a$ is the space
of holomorphic sections of the positive degree-$N$ bundle.  Riemann--Roch
on a genus-one curve and vanishing for the inverse negative-degree bundle
give $\dim\ker a=N$, for every flat twist.

The compact elliptic operator has discrete spectrum.  Repeated application
of $a$ to an eigenvector must terminate, because \eqref{eq:ladder} lowers
energy by $B$ while $H_B\ge B/2$.  It terminates in $\ker a$, so every
eigenvalue is one of \eqref{eq:levels}.  Conversely $(a^*)^n$ takes
$\ker a$ injectively onto the $n$th level; the lowering argument makes it
surjective.  Hence every multiplicity equals $N$.  Complex conjugation
replaces this by $|N|$ when $B<0$.
\end{proof}

\section{Finite magnetic translations}

Let $M=|N|$, $\sigma=\sgn N$, and $\zeta_M=\exp(2\pi\ii/M)$.  Fix the
positively oriented ordered $M$-division vectors
$(L_x/M,0)$ and $(0,L_y/M)$.  Their translations lift projectively to the
magnetic bundle.  After harmless scalar rephasings, call the corresponding
lifts $U,V$.

\begin{theorem}[Finite Heisenberg action]
On every Landau eigenspace one can choose an orthonormal basis
$e_0,\ldots,e_{M-1}$, with indices modulo $M$, such that
\begin{equation}
 Ue_j=e_{j+1},\qquad Ve_j=\zeta_M^{-\sigma j}e_j,
 \qquad U^M=V^M=I,\qquad UV=\zeta_M^\sigma VU.  \label{eq:clockshift}
\end{equation}
For $M>1$ this representation is irreducible; for $M=1$ it is the trivial
one-dimensional boundary case.  Flat holonomies change the lifts by
phases, not the commutator or the Landau data.
\end{theorem}

\begin{proof}
The bundle cocycle in Theorem~\ref{thm:flux} makes the commutator of the two
fixed-orientation lifted division translations equal to the exponential of
the flux through one $M$-by-$M$ cell.  Since
\[
 \frac{BA}{M^2}=\frac{2\pi N}{M^2}=\frac{2\pi\sigma}{M},
\]
this phase is $\zeta_M^\sigma$; in particular, negative flux cannot be
removed by scalar rephasing.  Rephasing does make both $M$th powers the identity.
They commute with $H_B$, hence act on each eigenspace.  Diagonalize $V$.
The commutator cycles its $M$ distinct characters, yielding
\eqref{eq:clockshift}.  Since the eigenspace has dimension $M$, the cycle is
a basis and any invariant subspace containing one character contains all.
The raising operator intertwines the action, so the same model occurs on
every level.
\end{proof}

\noindent\textit{Round-one certificate: round one adds the complete
Bochner Landau ladder, exact multiplicity, holonomy rigidity, and the
finite magnetic-translation representation.}
\fi

\ifnum\CRevisionRound>1
\section{Heat trace, spectral zeta, and determinant}

\begin{theorem}[Three exact spectral invariants]\label{thm:zeta}
For $\beta>0$ and initially for $\Re s>1$,
\begin{align}
 \Tr(\e^{-\beta H_B})
   &=M\frac{\e^{-\beta|B|/2}}{1-\e^{-\beta|B|}}, \label{eq:heat}\\
 \zeta_{H_B}(s)
   &=M|B|^{-s}\zeta(s,\tfrac12)
     =M|B|^{-s}(2^s-1)\zeta(s).                 \label{eq:zeta}
\end{align}
The meromorphic continuation is regular at zero and
\begin{equation}
 \zeta_{H_B}(0)=0,\qquad
 \zeta'_{H_B}(0)=-\frac{M}{2}\log2,\qquad
 \det_{\zeta}H_B=2^{M/2}.                       \label{eq:det}
\end{equation}
In particular, the zeta determinant is independent of $|B|$ once the flux
multiplicity $M$ is fixed.
\end{theorem}

\begin{proof}
Insert \eqref{eq:levels} with multiplicity $M$.  The heat sum is geometric,
giving \eqref{eq:heat}; the Dirichlet sum gives the first equality in
\eqref{eq:zeta}.  Separating the odd positive integers proves
$\zeta(s,1/2)=(2^s-1)\zeta(s)$ and supplies the continuation.  The standard
Hurwitz values
\[
 \zeta(0,a)=\tfrac12-a,\qquad
 \zeta'(0,a)=\log\Gamma(a)-\tfrac12\log(2\pi)
\]
give $\zeta(0,1/2)=0$ and $\zeta'(0,1/2)=-(\log2)/2$.  Differentiating the
factor $|B|^{-s}$ contributes nothing because the value at zero vanishes.
Finally use $\det_\zeta H_B=\exp[-\zeta'_{H_B}(0)]$.
\end{proof}

\paragraph{Interpretive boundary.}
Equation \eqref{eq:zeta} is a source spectral zeta.  This paper gives it no
dynamical-determinant, Euler-product, or target-function interpretation.

\section{Two exact revival times}

\begin{theorem}[Least scalar and identity revivals]\label{thm:revival}
The least positive time at which the propagator is a scalar is
\[
 T_{\rm scalar}=\frac{2\pi}{|B|}=T_{\rm cl},
 \qquad \e^{-\ii T_{\rm scalar}H_B}=-I.
\]
The least positive identity time is
\[
 T_{\rm identity}=\frac{4\pi}{|B|}=2T_{\rm cl},
 \qquad \e^{-\ii T_{\rm identity}H_B}=I.
\]
\end{theorem}

\begin{proof}
The ratio of phases on adjacent levels is $\exp(-\ii|B|t)$.  A scalar
propagator therefore requires and, by \eqref{eq:levels}, is implied by
$|B|t\in2\pi\ZZ$.  At the least positive such time the $n$th phase is
$\exp[-2\pi\ii(n+1/2)]=-1$.  At scalar times $2\pi k/|B|$, this common phase
is $(-1)^k$, so the least identity occurs at $k=2$.
\end{proof}

\section{Boundary atlas}

\begin{center}
\begin{tabular}{>{\raggedright\arraybackslash}p{0.20\linewidth}
                >{\raggedright\arraybackslash}p{0.72\linewidth}}
\toprule
face & exact status\\
\midrule
$E>0$, $B\ne0$ & Theorem~\ref{thm:classical}; all orbits have the common
least period and form a clean family.\\
$E=0$, $B\ne0$ & $p=0$ gives the zero-section of stationary classical
points.  Quantum mechanics still has positive ground energy $|B|/2$.\\
$N>0$ versus $N<0$ & Reversing $B$ reverses classical orientation and
complex-conjugates quantum chirality; energies and multiplicities use
$|B|$ and $|N|$, while fixed positive division vectors give conjugate
commutator phases $\zeta_M^{\pm1}$.\\
$|N|=1$ & Every Landau level is simple and the finite Heisenberg action is
one-dimensional; all spectral formulas remain valid.\\
$N\notin\ZZ$ & The classical twisted flow remains valid, but no global
Hermitian line bundle has the stipulated curvature.\\
$B=0$ & Classical motion is $q(t)=q_0+tp$ modulo the lattice and closes iff
$p_xt\in L_x\ZZ$ and $p_yt\in L_y\ZZ$ for some $t>0$.  For $p_xp_y\ne0$
this is equivalent to $p_yL_x/(p_xL_y)\in\mathbb Q$; nonzero axial motion
closes, while $p=0$ is stationary and has no positive least period.  A
nonaxial irrational normalized ratio gives a dense orbit.  The quantum
object becomes a flat-holonomy shifted torus Laplacian, not a zero-field
continuation of the Landau multiplicity theorem.\\
\bottomrule
\end{tabular}
\end{center}

\section{Route-A closure and reproducibility}

The strict tuple is
\[
 (\mathrm{A0\_FAIL},\mathrm{A1\_WEAK},
  \mathrm{A2\_FAIL},\mathrm{A3\_FAIL},
  \mathrm{A4\_NATURAL\_QUANTIZATION}),
\]
so the overall verdict is \texttt{ROUTE\_A\_REJECTED}.  Flux integrality is
a genuine source topological condition, but topology alone is not arithmetic
provenance and supplies no rational-prime data, prime-power repetition, or
logarithmic-prime clock; consequently A0 fails.
The classical returns are continuous clean families, and the exact Hurwitz
zeta is not a dynamical determinant.  No target continuation or zero
correspondence is asserted.  Route B is locked.

The companion artifact checks 256 exact classical cells, 16,512 Landau
labels, 4,160 signed clock--shift basis cells, 1,024 heat cells, 64 determinant
controls, 129 revival phases, and every listed boundary.  Producer,
independent checker, symbolic cross-check, byte replay, hostile mutation,
and PDF gates are separate.  These computations audit formula
transcription; the proofs above carry the infinite statements.  The scope
literal is \texttt{NO\_BAD\_EULER\_OR\_ROOT\_NUMBER}.

\noindent\textit{Round-two certificate: round two adds exact heat and
Hurwitz zeta closure, the zeta determinant, both least revival times, the
full boundary atlas, and the strict Route-A stop.}

\section*{Source and ownership note}
The torus magnetic-translation source is M.~H.~Al-Hashimi and
U.-J.~Wiese, \emph{Annals of Physics} 324 (2009), 343--360,
doi:10.1016/j.aop.2008.07.006, arXiv:0807.0630.  E.~Onofri,
\emph{International Journal of Theoretical Physics} 40 (2001), 537--549,
doi:10.1023/A:1004115827959, arXiv:quant-ph/0007055, treats the
finite-volume Landau-level subtlety.
I.~Burban and S.~Klevtsov, \emph{Communications in Mathematical Physics}
406 (2025), article 97, doi:10.1007/s00220-025-05267-9, provides modern
line-bundle and finite-Heisenberg context.  Its holomorphic ground-space
normalization is not substituted for the physical Bochner half-shift used
here.  Nearby repository owners cover the Penning trap, hyperbolic and
Grushin magnetic systems, the spherical Dirac monopole, and Harper lattice
dynamics.  This package claims only the displayed joint flat-torus closure,
not global literature novelty and not ownership of its standard ingredients.
\fi

\end{document}
